\section{实数连续性的等价描述}

在本书第一章我们讲述了实数基本定理，它断言实数系在戴德金意义下是连续的．我们还用它证明了单调有界原理，然后用单调有界原理证明了区间套定理，再用区间套定理证明了闭区间上连续函数的基本性质，由此已初步体会到实数基本定理的重要性．在本节我们将给出实数系连续性的其他等价描述．这些等价描述，虽然描写的都是实数连续性这同一件事，但由于出发的角度不同，在不同的情形下运用哪一个便有繁简之分，而且有些概念像下面要叙述的上、下确界，本身便有独立的意义，从而有广泛的应用，这可在函数可积性的讨论中窥见一斑．

我们知道，一个只有有限个元素的实数集，它总有最大数与最小数．但具有无穷个元素的实数集就不一定了．例如
\[
A = \left\{ \frac{n}{n+1} \;\middle|\; n = 1, 2, \cdots \right\},
\]
\[
B = \left\{ x \mid x \in \mathbf{Q}, x > 0 \text{ 且 } x^2 < 2 \right\},
\]
数集 $A$ 有最小数 $\frac{1}{2}$，没有最大数，但它有上界，无穷多个上界．上界中有一个最小的，这就是 $1$．数集 $B$ 既没有最大数，也没有最小数，但它有上界也有下界，而且有一个最小的上界 $\sqrt{2}$，有最大的下界 $0$．

\begin{definition}\label{def:9.1}
如果实数集合 $A$ 有上界，且上界中有最小数 $\beta$，则称 $\beta$ 为数集 $A$ 的上确界，记为
\[
\beta = \sup A;
\]
如果实数集合 $A$ 有下界，且下界中有最大的数 $\alpha$，则称 $\alpha$ 为数集 $A$ 的下确界，记为
\[
\alpha = \inf A.
\]
这里 $\sup$ 是 supremum 的缩写，而 $\inf$ 是 infimum 的缩写．
\end{definition}

由定义容易看出，$\beta$ 是数集 $A$ 的上确界，当且仅当
\begin{enumerate}[(i)]
\item 对任意 $x \in A$，有 $x \leqslant \beta$；
\item 对任意 $\varepsilon > 0$，存在 $x_{\varepsilon} \in A$，使得
\[
x_{\varepsilon} > \beta - \varepsilon.
\]
\end{enumerate}
事实上，条件~(i) 说的是，$\beta$ 是集合 $A$ 的一个上界；条件~(ii) 说的是，任何较 $\beta$ 小的数就不再是 $A$ 的上界了．因此，(i)、(ii) 合起来是 $\beta = \sup A$ 的充分必要条件．由此可见，如果一个数集有上确界存在，则上确界是唯一的．

由此也可见，如果集合 $A$ 有最大数，则此最大数便是集合 $A$ 的上确界．这时称这集合达到它的上确界．

类似地，$\alpha$ 是数集 $A$ 的下确界，当且仅当
\begin{enumerate}[(i)]
\item 对任意 $x \in A$，有 $\alpha \leqslant x$；
\item 对任意 $\varepsilon > 0$，存在 $x_{\varepsilon} \in A$，使得 $x_{\varepsilon} < \alpha + \varepsilon$．
\end{enumerate}

然而，有上界的数集是否一定有上确界呢？这在实数系内是成立的．

\begin{theorem}[确界定理]\label{thm:9.1}
在实数系 $\mathbb{R}$ 内，非空的有上（下）界的数集必有上（下）确界存在．
\end{theorem}

\begin{proof}
设 $X$ 是有上界的非空实数集，记 $B$ 为 $X$ 的全体上界组成的集合，$A = \mathbb{R} \setminus B$，则 $A \mid B$ 构成实数的一个分划．事实上，不空、不漏显然．我们来证“不乱”．任意 $a \in A$，$b \in B$，由 $a$ 不是 $X$ 的上界，知有 $x_0 \in X$，使得 $x_0 > a$，而由 $b \in B$ 知，$x_0 \leqslant b$，故 $a < b$．

由实数基本定理，分划 $A \mid B$ 确定唯一的实数 $r$，使得任意 $a \in A$，$b \in B$，有 $a \leqslant r \leqslant b$．我们要证明 $r = \sup X$．首先证明 $r$ 是 $X$ 的上界，用反证法，如果不然，则有 $x_0 \in X$，使得 $x_0 > r$，这时有 $a = \frac{x_0 + r}{2} \in A$，且 $a > r$，这是不可能的，因此 $r$ 是 $X$ 的上界．而由于任意 $b \in B$，$r \leqslant b$，这就表明了 $r$ 是 $X$ 的最小上界．下确界情形的证明是类似的，定理证完．
\end{proof}

下面我们来证明，定理~\ref{thm:9.1} 是和实数基本定理等价的．也就是说，我们刚才用实数基本定理证明了定理~\ref{thm:9.1}，现在反过来，可以用定理~\ref{thm:9.1} 证明实数基本定理．设给定 $\mathbb{R}$ 的一个分划 $A \mid B$．由于 $B$ 的任何一个数都是集合 $A$ 的上界，根据定理~\ref{thm:9.1} 知，$A$ 有上确界 $r$，显然对任意 $a \in A$ 有 $a \leqslant r$，而对任意 $b \in B$，由于 $b$ 是 $A$ 的上界，$r$ 是上确界，故有 $r \leqslant b$，实数基本定理证完．

我们已经证明了，定理~\ref{thm:9.1} 和实数基本定理是等价的，而实数基本定理说的正是实数系（按戴德金意义）是连续的，因此，定理~\ref{thm:9.1} 给出了实数连续性的一个等价描述．

实数连续性还有另外一个等价描述，那就是单调上升有上界的实数序列必有极限存在．我们在第三章中曾用实数基本定理证明过这一事实，下面用这个事实，反过来证明实数基本定理．

设给定实数 $\mathbb{R}$ 的一个分划 $A \mid B$．我们要用单调上升有上界数列必有极限的定理，证明存在实数 $r$，使得任意 $a \in A$，$b \in B$ 有 $a \leqslant r \leqslant b$．事实上，任取 $a_1 \in A$，$b_1 \in B$，用 $a_1$，$b_1$ 的中点 $\frac{a_1 + b_1}{2}$ 二等分 $[a_1, b_1]$．如果 $\frac{a_1 + b_1}{2} \in B$，则取 $a_2 = a_1$，$b_2 = \frac{a_1 + b_1}{2}$；如果 $\frac{a_1 + b_1}{2} \in A$，则取 $a_2 = \frac{a_1 + b_1}{2}$，$b_2 = b_1$．继续用 $\frac{a_2 + b_2}{2}$ 二等分 $[a_2, b_2]$．如果 $\frac{a_2 + b_2}{2} \in A$，则取 $a_3 = \frac{a_2 + b_2}{2}$，$b_3 = b_2$；如果 $\frac{a_2 + b_2}{2} \in B$，则取 $a_3 = a_2$，而 $b_3 = \frac{a_2 + b_2}{2}$．如此继续下去，便得到两串序列 $\{a_n\}$，$\{b_n\}$，其中 $a_n \in A$ 单调上升有上界（例如 $b_1$），$b_n \in B$ 单调下降有下界（例如 $a_1$）．并且 $b_n - a_n = \frac{b_1 - a_1}{2^n} \to 0$（$n \to \infty$）．由 $a_n$ 单调上升有上界知有 $r$ 存在，使得 $r = \lim\limits_{n \to \infty} a_n$．我们来证明，任意 $a \in A$，$b \in B$，有 $a \leqslant r \leqslant b$．这等价于证明，任意 $a < r$，$b > r$，有 $a \in A$，$b \in B$．事实上，任意 $a < r$，由于 $\lim\limits_{n \to \infty} a_n = r$，知存在 $n_0$，使 $a < a_{n_0} \in A$，故 $a \in A$，而对任意 $b > r$，由 $b_n = a_n + \frac{b_1 - a_1}{2^n}$ 知 $\lim\limits_{n \to \infty} b_n = r$，故存在 $n_1$ 使得 $b > b_{n_1} \in B$，从而 $b \in B$．这就证明了实数基本定理（读者可能也明白，我们在这里用了区间套的思想）．也就是说，现在我们已经证明了，实数基本定理同单调上升有上界的数列必有极限存在定理是等价的，即后者也给出了实数连续性的另一个等价描述．

把上面的推理总结一下，实数连续性有下述三种描述：
\begin{enumerate}[(1)]
\item 对实数系的每一个分划 $A \mid B$，存在唯一的实数 $r$，使得对任意 $a \in A$，$b \in B$，有 $a \leqslant r \leqslant b$．
\item 非空有上界的实数子集必有上确界存在．
\item 单调上升有上界的实数列必有极限存在．
\end{enumerate}
我们上面证明了 (1)$\Leftrightarrow$(2)，(1)$\Leftrightarrow$(3)．

上述的 (2) 与 (3)，对有理数系是不正确的．例如，前面举的
\[
A = \left\{ x \mid x \in \mathbf{Q}, x > 0 \text{ 且 } x^2 < 2 \right\}
\]
是有上界的有理数集，但它在有理数系 $\mathbf{Q}$ 中没有上确界（因为在实数系中它的上确界是 $\sqrt{2}$，不是有理数），这是由于有理数系在 $\sqrt{2}$ 处出现了空隙．类似地，有理数数列 $x_n = \left( 1 + \frac{1}{n} \right)^n$ 是单调上升有上界的数列，但它在有理数系中没有极限（因为在实数系中它的极限是 $e$，不是有理数），这是由于有理数系在 $e$ 的地方出现了空隙．这也就说明了 (2) 与 (3) 分别是实数连续性的等价描述的直观解释，不过这种直观远不如 (1) 那么强就是了．

值得指出的是，我们当然也可以直接证明 (2)$\Leftrightarrow$(3)，其中的 (2)$\Rightarrow$(3) 比较简单，我们把它写出来．

设 $\{x_n\}$ 是单调上升有上界的实数列．由 (2) 知 $r = \sup\{x_n\}$ 存在．且有 $x_n \leqslant r$．另外，任给 $\varepsilon > 0$，存在 $x_N > r - \varepsilon$，因此，当 $n > N$ 时有 $r - \varepsilon < x_N \leqslant x_n \leqslant r$，即 $|x_n - r| < \varepsilon$．这就证明了 $r = \lim\limits_{n \to \infty} x_n$，也就是证明了 (2)$\Rightarrow$(3)．

上述证明还说明，单调上升有上界的序列，其极限刚巧就是 $\{x_n\}$ 的上确界．

读者不妨自己写出 (3)$\Rightarrow$(2) 的证明．

\begin{xiti}
\item 求数列 $\{x_n\}$ 的上、下确界（若 $\{x_n\}$ 无上（下）界，则称 $+\infty (-\infty)$ 是 $\{x_n\}$ 的上（下）确界）：
\begin{ziti}
\item $x_n = 1 - \frac{1}{n}$；
\item $x_n = n \left[ 2 + (-2)^n \right]$；
\item $x_{2k} = k$，$x_{2k+1} = 1 + \frac{1}{k}$ ($k = 1, 2, 3, \cdots$)；
\item $x_n = \left[ 1 + (-1)^n \right] \frac{n+1}{n}$；
\item $x_n = \sqrt[n]{1 + 2^{n(-1)^n}}$；
\item $x_n = \frac{n-1}{n+1} \cos \frac{2n\pi}{3}$．
\end{ziti}
\item 设 $f(x)$ 在 $D$ 上定义，求证：
\begin{ziti}
\item $\sup\limits_{x \in D} \{ -f(x) \} = -\inf\limits_{x \in D} f(x)$；
\item $\inf\limits_{x \in D} \{ -f(x) \} = -\sup\limits_{x \in D} f(x)$．
\end{ziti}
\item 设 $\beta = \sup E$，且 $\beta \notin E$，试证自 $E$ 中可选取数列 $\{x_n\}$ 且 $x_n$ 互不相同，使 $\lim\limits_{n \to \infty} x_n = \beta$；又若 $\beta \in E$，则情形如何？
\item 试证收敛数列必有上确界和下确界，趋于 $+\infty$ 的数列必有下确界，趋于 $-\infty$ 的数列必有上确界．
\item 试分别举出满足下列条件的数列：
\begin{ziti}
\item 有上确界无下确界的数列；
\item 含有上确界但不含有下确界的数列；
\item 既含有上确界又含有下确界的数列；
\item 既不含有上确界又不含有下确界的数列，其中上、下确界都有限．
\end{ziti}
\item 设 $f(x)$ 在 $[a, b]$ 上有界，定义
\[
\omega_f[a, b] = \sup_{x \in [a, b]} f(x) - \inf_{x \in [a, b]} f(x),
\]
求证：
\[
\omega_f[a, b] = \sup_{x', x'' \in [a, b]} |f(x') - f(x'')|.
\]
\item 设 $f(x)$ 在 $x_0$ 附近有定义且有界，定义
\[
\omega_f(x_0) = \lim_{n \to +\infty} \omega_f\left[ x_0 - \frac{1}{n}, x_0 + \frac{1}{n} \right],
\]
求证：$f(x)$ 在 $x_0$ 连续的充分必要条件为 $\omega_f(x_0) = 0$．
\item 设 $f(x)$，$g(x)$ 在 $D$ 上有界，且大于 $0$，求证：
\begin{align*}
\inf_{x \in D} f(x) \cdot \inf_{x \in D} g(x) &\leqslant \inf_{x \in D} \{ f(x) g(x) \} \leqslant \inf_{x \in D} f(x) \cdot \sup_{x \in D} g(x) \\
&\leqslant \sup_{x \in D} \{ f(x) g(x) \} \leqslant \sup_{x \in D} f(x) \cdot \sup_{x \in D} g(x).
\end{align*}
\end{xiti}
